% =====================================================================
\chapter{Bilan critique}
\label{ch:bilan}
% =====================================================================

Un rapport de stage qui ne présente que des réussites n'est pas un rapport :
c'est une plaquette. Ce chapitre expose ce qui a été difficile, ce qui a mal
tourné, et ce que je ferais différemment.

\section{Ce qui a bien fonctionné}

\subsection{La remise en question du cadrage}

Le moment le plus productif du stage a été la session de recentrage du
25 août. Elle a consisté à admettre qu'un travail déjà réalisé — quatre
connecteurs, une couche de notation, un tableau de bord — n'apportait pas la
valeur attendue, et à le supprimer.

Cette décision a été facilitée par un critère unique et brutal :
\emph{l'acheteur pourrait-il obtenir ceci directement depuis la source ?} Un
critère unique est plus facile à appliquer qu'une liste de priorités, parce
qu'il ne laisse pas de place à la négociation.

\subsection{La confrontation au terrain métier}

Les échanges avec un actuaire et un professionnel de l'assurance ont
directement réécrit la direction produit. Le refus du score composite, la
préférence pour les variables brutes et l'exigence de provenance viennent de
ces conversations, pas d'une intuition.

La leçon est générale : \textbf{il est moins coûteux de se tromper auprès d'un
utilisateur que de se tromper tout seul pendant deux mois}. Ces échanges ont
coûté quelques heures et ont évité de construire un produit que personne
n'aurait voulu.

\subsection{La discipline documentaire}

Le triptyque constitution / spécification / plan a permis de stabiliser un
projet dont une part importante du code a été écrite en collaboration avec des
agents de programmation. Sans ce cadre, le projet aurait très probablement
dérivé : chaque session d'agent aurait pu défaire silencieusement une décision
antérieure.

\section{Ce qui a mal fonctionné}

\subsection{Les bugs sémantiques répétés}

Trois fois de suite, un comportement \emph{fonctionnel} a dû être corrigé parce
qu'il était \emph{faux} :

\begin{enumerate}[leftmargin=1.6em]
  \item les clés de scénario ne correspondaient pas ;
  \item un seuil unique confondait dommages et affichage ;
  \item le niveau d'eau dépendait de la météo du jour.
\end{enumerate}

Dans les trois cas, le code s'exécutait sans erreur. Le fait que ces bugs
soient apparus \emph{après} la mise en service montre une faiblesse de méthode :
les tests vérifiaient que le code ne plantait pas, pas qu'il disait vrai. Une
partie de la réponse a été la suite de tests de scénarios
(\fichier{scenarioLoop.test.ts}), qui verrouille la chaîne complète de la clé
jusqu'à l'eau dessinée. Mais cette suite est arrivée \emph{après} les bugs.

Ce que je retiens : \textbf{pour un produit de données, un test qui vérifie
l'absence d'exception est presque inutile}. Ce qui compte est un test qui
vérifie une valeur attendue sur une entrée connue.

\subsection{La dérive initiale}

La phase 1 a produit du code qui a dû être supprimé. Avec le recul, le
problème n'était pas le travail lui-même — explorer quatre sources était
nécessaire pour comprendre ce qui avait de la valeur — mais l'\emph{absence de
critère explicite} pendant cette exploration. Il a fallu construire le critère
d'abord, puis mesurer le travail à l'aune de ce critère.

Si c'était à refaire, j'écrirais la question \og l'acheteur pourrait-il obtenir
ceci depuis la source ? \fg{} dès la première semaine, avant d'écrire la
première ligne de connecteur.

\subsection{L'instabilité du déploiement}

Le déploiement a consommé une part disproportionnée du temps. Les messages
d'erreur de la plateforme désignaient systématiquement le symptôme, jamais la
cause. Le message \og le motif ne correspond à aucune fonction \fg{} a coûté
plusieurs itérations alors que la cause réelle — la visibilité du nom racine
dans le module — n'était mentionnée nulle part.

La leçon : \textbf{sur une plateforme d'hébergement, il faut diagnostiquer par
l'expérience contrôlée, pas par le message d'erreur}. Réduire la surface,
changer une variable à la fois, observer.

\subsection{Le périmètre du mode VFX}

Le mode VFX a été spécifié, encadré, partiellement implémenté — et non terminé.
Trois tâches sur dix restent ouvertes. C'est un périmètre trop ambitieux pour
la durée disponible, et la conséquence est un demi-livrable.

Le point positif est que la partie livrée est la partie \emph{gouvernance} :
l'amendement, l'étiquetage, le bus d'entrée et les garanties. C'est la partie
qui protège le produit. Le rendu 3D lui-même est la partie amusante, et c'est
celle qui manque.

\section{Ce que je ferais différemment}

\begin{table}[H]
\centering
\small
\caption{Choix passés et alternatives envisagées.}
\label{tab:retro}
\begin{tabularx}{\textwidth}{@{}X X@{}}
\toprule
\textbf{Ce qui a été fait} & \textbf{Ce que je ferais autrement} \\
\midrule
Explorer quatre sources avant de définir le critère de valeur &
  Écrire le critère de valeur dès le premier jour, puis explorer \\
\addlinespace
Écrire les tests de scénarios après la mise en service &
  Écrire un test de valeur attendue (clé $\rightarrow$ pic $\rightarrow$
  profondeur) dès l'implémentation du moteur \\
\addlinespace
Traiter le déploiement au fil de l'eau &
  Déployer une version minimale dès la première semaine, pour découvrir
  l'infrastructure tôt plutôt que tard \\
\addlinespace
Spécifier dix tâches VFX &
  Spécifier trois tâches VFX (bascule, étiquetage, garanties) et garder le
  rendu pour une phase ultérieure \\
\addlinespace
Garder la calibration de latence pour la fin &
  Mesurer dès que le pipeline est stable, pour que le budget soit un
  engagement tenu et non une hypothèse \\
\bottomrule
\end{tabularx}
\end{table}

\section{Compétences acquises}

\begin{table}[H]
\centering
\small
\caption{Compétences mobilisées et développées.}
\label{tab:competences}
\begin{tabularx}{\textwidth}{@{}l X@{}}
\toprule
\textbf{Domaine} & \textbf{Acquis} \\
\midrule
Systèmes d'information géographique &
  WFS et GML 3.2, test point-in-polygon par lancer de rayons, proximité
  géométrique, projections et ordre des axes, boîtes englobantes \\
\addlinespace
Ingénierie backend &
  FastAPI, Pydantic v2, asynchrone avec \code{httpx}, conception de contrat de
  données, gestion de cache et de quotas \\
\addlinespace
Ingénierie frontend &
  React et TypeScript, Mapbox GL JS (extrusion, couches personnalisées,
  terrain), animation synchronisée à un état applicatif \\
\addlinespace
Qualité logicielle &
  Tests hors-ligne avec fixtures, tests de dégradation, tests de propriété,
  portes de recherche CI \\
\addlinespace
Méthodologie &
  Développement piloté par les spécifications, gouvernance par document,
  collaboration structurée avec des agents de programmation \\
\addlinespace
Déploiement &
  Architecture serverless mono-domaine, réécritures et proxies, diagnostic
  d'incidents de plateforme \\
\addlinespace
Domaine métier &
  Souscription, Solvabilité II, régime CatNat, non-stationnarité du risque
  climatique, distinction entre donnée et jugement \\
\bottomrule
\end{tabularx}
\end{table}

\section{Appréciation personnelle}

Ce stage m'a surtout appris qu'un projet technique se joue moins sur la
difficulté du code que sur la justesse des décisions. Écrire un parseur GML
est un exercice ; décider qu'il ne faut \emph{pas} produire de score est une
décision. La seconde a demandé plus de courage et a eu plus d'impact.

J'en retiens une posture : chercher la formulation la plus honnête possible du
problème, et s'y tenir même quand elle impose de supprimer du travail. Le
recentrage du 25 août a été inconfortable sur le moment — il fallait admettre
que plusieurs semaines de production ne servaient pas le produit — mais c'est
ce qui a donné au livrable sa cohérence.

Le projet n'est pas terminé. Il est, en revanche, défendable : chaque choix est
documenté, chaque fait porte sa source, et chaque limite connue est écrite
plutôt que cachée. C'est probablement la définition d'un travail honnête.
